% ================================================================ conclusion
% Unnumbered chapter, per the OBE 2.1 template ToC.
\chapter*{Conclusion}
\addcontentsline{toc}{chapter}{Conclusion}

This thesis addressed the inadequacy of binary deepfake detection in a media
landscape where the visual and acoustic streams of a clip can be manipulated
independently. The stated objective---to detect, per modality, whether the
video and the voice of a clip are real or AI-generated, to attribute the clip
to one of four authenticity quadrants, and to localize manipulated segments in
time---was pursued through a designed artifact combining DAVID-Net, a
disentangled cross-modal transformer, with Quadrant-Aware Contrastive
Pretraining, a generator-free pretraining strategy that constructs all four
quadrants from pristine data alone.

The data collection and analysis process relied on established benchmark
corpora (FakeAVCeleb, AV-Deepfake1M, LAV-DF) with cross-dataset validation on
unseen corpora (DFDC, KoDF), unified under a single reproducible manifest
schema with subject-disjoint and leave-one-generator-out splits.

The principal findings are threefold (all values are draft placeholders,
shown in \textcolor{blue}{blue}, pending the final experimental runs).
First, a single network can attribute manipulation per modality at the level
of specialized unimodal detectors---\dummy{0.986} video AUC and
\dummy{0.991} audio AUC with \dummy{94.6\%} quadrant accuracy on
FakeAVCeleb---demonstrating that binary detection discards information it
never needed to discard. Second, the generator-free QACP pretraining is the
decisive factor for generalization: it contributes \dummy{+0.065} AUC on
unseen datasets and lifts leave-one-generator-out AUC from \dummy{0.849} to
\dummy{0.920}, while in-domain accuracy is nearly unchanged---evidence that
what matters for real-world deployment is learned from the structure of
manipulation itself rather than from generator fingerprints. Third, the
ablation study shows that the disentanglement of authenticity from
consistency evidence, enforced by the MISMATCH pseudo-class, is what allows
these two capabilities to coexist in one model. The complete system is
released as open-source software and deployed as a public web service.

\textbf{Limitations.} The framework targets single-subject talking-head
content; multi-person scenes, overlapping speech, and non-speech audio are out
of scope. Detection outputs are probabilistic decision support, not proof of
authenticity. Empirically, calibration degrades under dataset shift
(ECE \dummy{0.064} on DFDC without recalibration), performance in the
combined low-light, low-resolution regime remains the weakest operating
point (video AUC \dummy{0.883}), and residual fairness gaps of up to
\dummy{0.024} AUC across apparent skin-tone subgroups persist after
balanced sampling; these are documented in the deployed model card.

\textbf{Future work.} Future researchers can extend this study toward fully
generated text-to-video content without a real source subject; continual
adaptation to emerging generators using the QACP construction as a
self-refreshing training signal; multi-person and non-frontal scenarios;
integration with provenance standards such as C2PA content credentials; and
on-device, real-time variants of the detector.
